% Chapter 5: Modifying the interaction terms
\chapter{Modifying the interaction terms}\label{Ch5:ModifyingInteraction}

\section{Constrained QND-like Protocol}

\qquad The full QND-like protocol requires an interaction Hamiltonian of the form $H_S(t) \otimes X_M$. For a QUBO problem Hamiltonian $H_P$, this necessitates simultaneous coupling to both the problem terms and the transverse driver $H_0$. Specifically, if $X_M = \sigma_z$, this implies engineering three-body terms like $\sigma_x \sigma_z$ (from the transverse field) alongside problem-based coupling. A more experimentally feasible "constrained" protocol restricts the coupling solely to the problem Hamiltonian $H_f$

\begin{equation}
    H_{\text{int}}(t) = f(t) H_f \otimes X_M.
\end{equation}

\qquad This approach eliminates the difficult $\sigma_x \sigma_z$ terms, requiring only variable-coupler interactions ($\sigma_z \sigma_z$ and $\sigma_z \sigma_z \sigma_z$). However, omitting the driver term means the interaction no longer commutes with the total Hamiltonian, $[H_{\text{int}}(t), H_S(t)] \neq 0$, thereby introducing relaxation effects alongside the intended energy rescaling.

\qquad Numerical results for this constrained protocol (Fig. \ref{fig:constrained_results_b}) reveal a saturation in performance. Unlike the ideal case, where fidelity improves linearly with interaction strength, the constrained protocol's fidelity plateaus near $x_0 \approx 2.0$. This limitation arises because stronger coupling increases the rate of deleterious relaxation transitions, eventually competing with the advantages of the widened energy gap. Thus, while the constrained protocol offers a speedup, it is bounded and cannot be indefinitely enhanced by increasing interaction strength.

\begin{figure}[htbp]
    \centering
    \begin{subfigure}[t]{0.34\textwidth}
        \centering
        \refstepcounter{subfigure}\makebox[\textwidth][l]{\hspace{-2mm}(\thesubfigure)}\label{fig:constrained_results_a}
        \vspace{-1.2ex}
        \includegraphics[width=\textwidth]{Dephasing/Dephasing5a.png}
    \end{subfigure}%
    \hfill
    \begin{subfigure}[t]{0.63\textwidth}
        \centering
        \refstepcounter{subfigure}\makebox[\textwidth][l]{\hspace{-2mm}(\thesubfigure)}\label{fig:constrained_results_b}
        \vspace{-1.2ex}
        \includegraphics[width=\textwidth]{Dephasing/Dephasing5bc.png}
    \end{subfigure}
    \caption[TTS speedup and fidelity for constrained QND protocol]{(a) Time-to-Solution (TTS) speedup ratio versus number of qubits $N$ for full and constrained QND protocols ($x_0=2$). The dashed line marks the theoretical limit $1/(1+x_0)$. (b) Average fidelity of the constrained protocol as a function of interaction strength $x_0$, showing the performance plateau due to non-commuting dynamics.}
    \label{fig:constrained_results}
\end{figure}

%----------------------------------------------------------

\section{Time-to-Solution Analysis}

\qquad To assess the practical utility of the engineered dephasing protocol, we employ the Time-to-Solution (TTS) metric and investigate a constrained implementation that relaxes the rigorous hardware requirements of the full QND interaction.

\qquad The Time-to-Solution (TTS) serves as a robust benchmark for quantum annealing, balancing the trade-off between the duration of a single anneal and the probability of reaching the ground state. Given that verifying a candidate solution is typically efficient, a probabilistic algorithm is effective provided it yields the correct result with a finite probability $p_{\text{single}}$. The TTS, denoted $T_p$, quantifies the expected total time required to achieve a solution with a target confidence level $p$ (typically set to 0.95)

\begin{equation}
    T_p = \min_T \frac{T \log(1-p)}{\log(1-p_{\text{single}}(T))}.
\end{equation}

\qquad Based on the energy rescaling derived in previous sections, the success probability for the QND-like protocol scales as $p_{\text{single}}^{\text{QND}}(T) = p_{\text{single}}^{\text{coh}}((1+m)T)$, where $m$ is the meter eigenvalue. Consequently, the ratio of the TTS for the QND-modified system to the standard coherent system is expected to perform as

\begin{equation}
    \frac{T_p^{\text{QND}}}{T_p^{\text{coh}}} = \frac{1}{1+m}.
\end{equation}

\qquad Numerical verification of this scaling is presented in Fig. \ref{fig:constrained_results_a} for random Ising instances. With a coupling strength of $x_0 = 2.0$, the observed speedup ratio converges to the theoretical prediction of $(1+x_0)^{-1} \approx 0.33$ for system sizes $N > 3$. Minor deviations in smaller systems act as artifacts of finite sampling within the minimization procedure over annealing time $T$.

%----------------------------------------------------------

\section{Interaction Graph Geometry}\label{sec:geometry}

\qquad The analysis presented in the original paper\cite{DephasingPaper} considers random Ising Hamiltonians with all-to-all connectivity, where every pair of qubits is coupled. However, many physical quantum annealing architectures impose constraints on qubit connectivity, limiting couplings to specific graph topologies. We extend the investigation by studying how the geometry of the interaction graph in the problem Hamiltonian $H_f$ affects the performance of both the QND-like and constrained QND-like protocols. Four representative geometries are considered: the \textit{chain}, \textit{ring}, \textit{star}, and \textit{wheel} graphs.

\subsection{Graph Definitions}

\qquad The four interaction geometries studied are defined as follows. For a system of $N$ qubits labelled $0, 1, \ldots, N-1$:
\begin{itemize}
    \item \textbf{Chain}: Sequential nearest-neighbour couplings, with edges $(i, i+1)$ for $i = 0, \ldots, N-2$. This yields $N-1$ edges.
    \item \textbf{Ring}: A chain with the addition of a periodic boundary edge $(0, N-1)$, giving $N$ edges in total.
    \item \textbf{Star}: A central hub qubit (qubit $0$) coupled to all others, with edges $(0, j)$ for $j = 1, \ldots, N-1$. This also yields $N-1$ edges, but with highly asymmetric connectivity.
    \item \textbf{Wheel}: A star graph supplemented by a ring among the outer qubits, giving $2(N-1)$ edges. Qubit $0$ connects to all others, and qubits $1, \ldots, N-1$ form a closed ring.
\end{itemize}
In each case, the coupling strengths $J_{ij}$ and local fields $h_i$ are drawn uniformly from $[0, 1]$ as in the original study. The Ising problem Hamiltonian retains the form of Eq. (\ref{IsingProblemEquation}), but with $J_{ij} = 0$ for qubit pairs not connected by an edge in the chosen geometry. The all-to-all case studied in the paper corresponds to $N(N-1)/2$ edges and serves as a baseline for comparison.

\subsection{QND-like Protocol: Geometry Independence}

\qquad For the full QND-like protocol, where $H_{\text{int}}(t) = H_S(t) \otimes X_M$, the interaction commutes with the system Hamiltonian at all times by construction. The energy rescaling mechanism is therefore exact: every eigenvalue of $H_S(t)$ is uniformly scaled by the factor $(1 + x_0)$, regardless of the underlying graph topology. As a result, the TTS speedup ratio $\mathcal{T}/\mathcal{T}^{\text{coh}}$ converges to the theoretical prediction of $1/(1+x_0)$ for all four geometries studied, consistent with the all-to-all case. The averaged QND-like TTS ratios across the chain, ring, star, and wheel geometries closely track the values obtained from the random all-to-all graph in the original study (Fig.~\ref{fig:geometry_tts}). This confirms that the QND-like speedup is a universal property of the protocol, independent of the problem graph structure, as the commuting interaction guarantees pure energy rescaling without introducing geometry-dependent relaxation pathways.

\subsection{Constrained Protocol: Geometry Dependence}\label{subsec:constrained_geometry}

\qquad In the constrained QND-like protocol, the situation is qualitatively different. Here, the interaction $H_{\text{int}}(t) = f(t) H_f \otimes X_M$ couples only the problem Hamiltonian to the meter, and this interaction does not commute with the full system Hamiltonian $H_S(t)$. The non-vanishing commutator $[f(t)H_f, (1-f(t))H_0] \neq 0$ introduces relaxation transitions alongside the intended energy rescaling, and the severity of these unwanted transitions depends on the structure of $H_f$.

\qquad Numerical results for the TTS speedup ratio as a function of system size $N$ are presented in Fig.~\ref{fig:geometry_tts}. The chain and ring geometries exhibit favourable performance, with TTS ratios decreasing toward $\sim 0.4$ as $N$ increases, indicating a meaningful speedup that improves with system size. The wheel geometry shows intermediate performance, stabilising around $\sim 0.5$. In contrast, the star geometry displays markedly degraded performance: the TTS ratio increases from approximately $0.79$ at $N=3$ to $0.94$ at $N=8$, approaching unity and indicating a near-complete loss of speedup for larger systems. This stands in direct contrast to the all-to-all constrained results from the original paper, which maintain a roughly constant TTS ratio around $0.6$--$0.65$ for $N \geq 5$.

\qquad The poor scaling of the star geometry under the constrained protocol can be attributed to two interrelated structural factors:

\begin{enumerate}
    \item \textbf{Asymmetric coupling concentration.} In the star graph, all $N-1$ edges pass through the central hub qubit. As $N$ grows, additional leaf qubits are introduced with no direct mutual coupling in $H_f$. The energy rescaling provided by the constrained interaction $f(t) H_f \otimes X_M$ therefore only amplifies gaps associated with transitions involving the hub qubit. Transitions between leaf qubits, which are coupled indirectly through the transverse field $H_0$ but not through $H_f$, receive no protection from the meter coupling.

    \item \textbf{Growing mismatch between $H_f$ and $H_S(t)$.} The transverse field $H_0 = \sum_i \sigma_x^{(i)}$ acts uniformly on all $N$ qubits, whereas the star-topology $H_f$ couples only through the central node. As the system size increases, the non-commuting component $[f(t)H_f, (1-f(t))H_0]$ grows relative to the beneficial energy rescaling, because the transverse terms on the leaf qubits remain entirely uncompensated. The consequent relaxation rate increases in tandem with the coupling strength, consistent with the saturation mechanism identified for the general constrained protocol: the unwanted relaxation becomes stronger alongside the desired energy rescaling effects\cite{DephasingPaper}.
\end{enumerate}

\qquad In contrast, the chain and ring topologies distribute edges across all qubits, ensuring that the constrained coupling rescales energy gaps more uniformly throughout the spectrum. The ring's closed boundary condition provides an additional edge, slightly improving homogeneity relative to the chain. The wheel geometry benefits from its dual structure, the star core provides strong hub-mediated coupling while the outer ring ensures direct leaf-to-leaf interactions, resulting in performance that interpolates between the star and the more distributed topologies.

\begin{figure}[htbp]
    \centering
    \includegraphics[width=0.75\textwidth]{InteractionTerms/GeometryTTS.png}
    \caption[TTS speedup ratio for different interaction graph geometries]{Time-to-Solution speedup ratio $\mathcal{T}/\mathcal{T}^{\mathrm{coh}}$ versus system size $N$ for different interaction graph geometries. The QND-like protocol (squares) achieves the theoretical prediction $1/(1+x_0)$ (green dashed line) regardless of geometry. The constrained protocol (circles) shows strong geometry dependence: chain and ring topologies maintain a substantial speedup, while the star geometry degrades toward no speedup as $N$ increases. Dashed curves with reduced opacity denote the all-to-all random results from the original study. All data are averaged over 100 random Ising instances with $x_0 = 2.0$.}
    \label{fig:geometry_tts}
\end{figure}

\qquad These results highlight a practical consideration for implementing constrained QND-like protocols on hardware with restricted connectivity: the topology of the interaction graph plays a critical role in determining the achievable speedup. Architectures with more uniformly distributed couplings, such as chains or rings, preserve significantly more of the speedup afforded by the energy rescaling mechanism. Conversely, highly centralised topologies like the star graph are ill-suited for the constrained protocol, as the hub bottleneck creates an increasingly unfavourable balance between energy rescaling and relaxation as the system scales.

\newpage